\chapter*{全书架构设计}
\addcontentsline{toc}{chapter}{全书架构设计}

本书按百万字级教材规划，采用三卷、十二个大章、若干附录的结构。每个大章对应一个完整专题，原先按课程周次拆开的细小章节已经合并为大章内部的节和小节；后续写作会在这些节内继续加厚，而不是继续增加大量短章。

\section*{写作目标}

本书最终目标是覆盖：

\begin{itemize}
  \item C/C++ 程序如何变成机器执行。
  \item x86-64 汇编、ABI、ELF、链接和加载。
  \item 编译器 IR、优化、自动向量化和未定义行为。
  \item 现代 CPU 微架构、cache、TLB、分支、乱序执行。
  \item 可信性能测量、benchmark、PMU 和 profiling。
  \item AVX2/FMA、AVX-VNNI、int8 量化。
  \item GEMM、packing、microkernel、MiniBLAS。
  \item softmax、layernorm、RMSNorm、GELU、attention、operator fusion。
  \item 线程、NUMA、ISA dispatch、生产库阅读。
\end{itemize}

\section*{十二个大章}

\begin{enumerate}
  \item 程序、机器与 C++ 底层地基。
  \item 二进制、链接、进程与运行时系统。
  \item x86-64 汇编、ABI 与 C++ 二进制接口。
  \item 可信性能测量、Benchmark 工程与 Lab 00。
  \item 编译器、LLVM IR 与优化。
  \item 现代 CPU 微架构性能模型。
  \item 内存层级、数据布局、TLB 与 Roofline。
  \item SIMD、AVX2/FMA 与低精度计算。
  \item HPC Kernel：从循环优化到 MiniBLAS。
  \item AI CPU 算子优化。
  \item 性能工具、并行、NUMA 与生产库阅读。
  \item Capstone：大师级性能工程项目。
\end{enumerate}

\section*{章节深度标准}

每个大章按三层写：

\begin{enumerate}
  \item 必须掌握：刚学完 C/C++ 的读者能跟上。
  \item 工程理解：能用于真实代码分析和优化。
  \item 深入方向：连接体系结构、OS、编译器或生产库更深内容。
\end{enumerate}

\section*{交付物标准}

每个大章最终应包含：

\begin{itemize}
  \item 不少于一个贯穿例子。
  \item 不少于三个实验。
  \item 不少于三组习题。
  \item 一个报告型大作业。
  \item 一组常见误区。
  \item 一组验收标准。
\end{itemize}
